% !TeX root = main_fps.tex
\section{Introduction}

Particle-based simulation methodologies represent physical systems as collections of discrete interacting particles. 
The evolution of these systems is governed by interaction rules that determine how particles influence one another through
 time. In existing particle methods, these interactions are commonly derived from constitutive relationships, equations of 
state, potential-energy functions, collision-response models, or other phenomenon-specific formulations. Although the 
mathematical forms of these models differ substantially, they generally share a common objective: the generation of 
forces that are subsequently integrated to produce particle motion and the evolving state of the system.


The assumptions used to define particle interactions play a central role in determining the behavior of a particle methodology.
Interaction models influence not only the forces generated within the system, but also the physical phenomena that may emerge, 
the numerical properties of the simulation, and the computational cost required to evaluate particle interactions. 
Consequently, the choice of interaction model often reflects a balance between physical representation, numerical stability, 
computational efficiency, and the specific phenomena the methodology is intended to reproduce.

Existing particle methodologies frequently employ interaction models that are tailored to particular classes of physical phenomena. 
Constitutive relationships are commonly introduced to represent material response, equations of state are used to model compressibility 
and pressure effects, and specialized formulations are often developed to reproduce specific behaviors such as turbulence, fracture,
 phase change, or collision response. While these approaches have demonstrated success across many application domains, they typically 
rely upon assumptions that are specific to the phenomena being modeled. Consequently, the development of a particle methodology often
 becomes closely coupled to the physical behaviors it is intended to reproduce.

The present work explores an alternative approach to particle dynamics through the Field Particle System (FPS). 
Rather than beginning with interaction models tailored to specific physical phenomena, FPS investigates a framework 
in which particle behavior emerges from geometric relationships between interacting particles. Within the methodology, 
geometry serves as the primary source of interaction, providing the information required to generate forces, evolve particle motion, 
and determine the future state of the system. The objective is not to reproduce a particular phenomenon through specialized 
governing models, but to establish a common interaction framework from which complex system behavior may emerge through repeated 
local particle interactions.

FPS is founded upon a set of principles that collectively define how particle interactions are represented and evaluated. 
Interactions are derived from geometric relationships between particles, contact is represented through interpenetrating overlap 
states, and particle response is determined from the current geometric configuration of the system. The methodology employs a 
source-owned execution model in which particles update only their own state, while energy is treated as a diagnostic quantity 
rather than a governing variable.Together, these principles establish an instantaneous, geometry-driven framework in which 
particle response is determined from the current geometric configuration of the system and state evolution is governed through 
source-owned interactions. The combination of instantaneous response and source-owned state ownership naturally supports parallel 
execution while allowing particle motion to emerge through local interactions and the continuous evolution of system geometry.

This work presents the Field Particle System (FPS), a particle-based simulation methodology founded upon geometric interactions, 
interpenetrating contacts, instantaneous dynamics, energyless state evolution, and source-owned execution. The paper defines the 
governing principles of the framework, describes the conversion of geometry into motion through local particle interactions, and 
details the numerical procedures used to evolve the system through time. The methodology is demonstrated through a series of 
simulation examples, including a three-dimensional cathedral environment containing interacting air, dust, and light particles, 
a three-dimensional incense simulation exhibiting the development of complex flow structures, and a two-dimensional converging-diverging 
nozzle simulation. Together, these demonstrations illustrate the ability of a common geometric interaction framework to produce a
 diverse range of particle behaviors and flow phenomena.


